\documentclass[12pt, a4paper]{article}

\usepackage[margin=1in]{geometry}
\usepackage{enumitem}
\usepackage[hidelinks]{hyperref}
\usepackage{titlesec}
\usepackage{setspace}
\usepackage{parskip}
\usepackage{amsmath}
\usepackage{import}
%\usepackage[scaled]{helvet}
\usepackage{fontspec}
\setmainfont{Times New Roman}

\usepackage{xcolor}

% Use sans serif font
%\renewcommand{\familydefault}{\sfdefault}

% Formatting sections
\titleformat{\section}{\large\bfseries}{}{0em}{}[\titlerule]
\titlespacing{\section}{0pt}{1.5em}{0.3em} % {left}{before}{after}

\definecolor{oxfordblue}{HTML}{1F3A60}


% Main entry macro
\newcommand{\entry}[3]{%
  \noindent
  \begin{tabular*}{\textwidth}{@{\extracolsep{\fill}} p{0.8\textwidth} r}
    #1 & {\small #2} \\
    \emph{#3} &
  \end{tabular*}%
  \vspace{0.3em}
}

\newcommand{\entryitem}[1]{\item[] {\small #1}}

% For the teaching subsection
\newcommand{\teachingentry}[2]{%
  \noindent
  \begin{tabular*}{\textwidth}{@{\extracolsep{\fill}} l r}
    \emph{#1} & { #2} \\
  \end{tabular*}
  \vspace{0.05em}%
}

% Entry without bold
\newcommand{\regentry}[2]{%
  \noindent
  \begin{tabular*}{\textwidth}{@{\extracolsep{\fill}} p{0.8\textwidth} r}
    #1 & { #2} \\
  \end{tabular*}%
  \vspace{0.3em}%
}


\setlength{\parindent}{0pt}
\renewcommand{\labelitemi}{--}

\begin{document}

% NAME AND CONTACT INFO
{\LARGE \textbf{Henry Zhang}} \\[0.2cm]
Email: {\color{oxfordblue}\href{mailto:henryzhang@cuhk.edu.hk}{henryzhang@cuhk.edu.hk}} \\
%Website: {\color{oxfordblue}\href{https://henryhzhang.com}{henryhzhang.com}} \\
%Phone: +852 3943 6454
\vspace{-1.5em}

% RESEARCH INTERESTS
\section*{Research Interests}
Corporate Finance, Macro-Finance, Financial Intermediation, Environmental Finance

% EMPLOYMENT
\section*{Employment}
\entry{CUHK Business School}{2025–Now}{Assistant Professor of Finance}

% EDUCATION
\vspace{-5mm}
\section*{Education}
\entry{Massachusetts Institute of Technology (MIT)}{2019–2025}{Ph.D. in Economics}
\entry{Swarthmore College}{2013–2017}{B.A. with highest honors in Mathematics and Economics}

% WORKING PAPERS
\vspace{-5mm}
\section*{Working Papers}
{\color{oxfordblue}
\noindent \href{https://henryhzhang.com/files/Factoring.pdf}{Corporate Liquidity Supply from Non-Bank Intermediaries and the Real Effects of Factoring }
}
(with Victor Orestes and Thiago Silva) 
%\\
% We show that firms experience large increases in sales and purchases after receiving cheaper liquidity. We focus on factoring, defined as the supplier-initiated sale of receivables. In Brazil, receivables funds (FIDCs) securitize receivables for institutional investors. By assembling a novel transaction-level dataset of factoring with other credit operations for all registered firms and FIDCs, we construct a shift-share instrument for factoring financing supply based on FIDC flows. We then use a novel combination of electronic payments, trade credit, and employer-employee matched data to estimate the impacts. A flow-induced increase in receivables demand reduces firms' factoring interest rate. In response, firms demand more permanent labor and less temporary labor. In our model, these effects arise from factoring's purpose of reducing cash inflow volatility, helping firms match inflows to outflows, which firms otherwise achieve at an efficiency cost through substitution across labor types. Using our model, we estimate that an aggregate decrease in the economy-wide factoring spread by 1 percentage point leads to 0.3 to 0.5 percentage point increases in aggregate output and wages. \\[0.2em]

{\color{oxfordblue}
\noindent \href{https://henryhzhang.com/files/Frost_Shock.pdf}{Volatility and Under-Insurance with Limited Pledgeability: Evidence from a Frost Shock}
} 
(with Victor Orestes and Thiago Silva) 
%\\
% We use transaction-level data on payments, credit, and insurance to examine how Brazilian farmers responded to the severe frost of July 2021, a shock that affected coffee, a perennial crop whose plants are a major component of farm value. The frost shock reduced both output and the pledgeable value of farmers' collateral. We find that insured farmers increased investment in the years following the shock, while uninsured farmers reduced investment and borrowing. We show how this pattern is consistent with models of imperfect pledgeability of a firm's collateral, where constrained firms neither insure (ex-ante) nor fully recover from a shock (ex-post). Limited commitment endogenously generates under-insurance through the combination of upfront payment of the insurance premium with the tightening of borrowing constraints post-shock due to the decrease in total collateral. We discuss two equilibrium implications of this mechanism: the inefficacy of emergency credit lines in targeting liquidity constrained firms and the amplification of output volatility from the rising risk of extreme weather shocks.\\[0.2em]


% RESEARCH IN PROGRESS
\section*{Research in Progress}
\noindent Aggregate Impacts of Command-and-Control Environmental Policy: Evidence from Court-Ordered Mining Bans in India (with Ananya Kotia and Utkarsh Saxena) %\\
% We estimate the aggregate impacts of court-ordered iron ore mining bans in India and consider the counterfactual welfare gains from an alternative policy to the ban. The local sectoral ban is a command-and-control (CAC) policy that is commonly applied to natural resource settings, usually when the regulator has a signal of widespread non-compliance. The Supreme Court of India imposed bans on iron ore mining and outbound iron ore trade in two states in response to reports that mines operated under fake environmental permits and underpaid mining royalties. Using firm-level industrial survey data, mine-level output data, and bilateral mine-to-firm auction data, we decompose the bans' effects into trade, production networks, and local labor demand channels. Our results indicate persistent declines in employment, capital stock, and borrowing by iron-consuming plants, despite the temporary duration of the ban. These findings highlight the economic spillovers caused by CAC policies, especially in industries that are upstream in the supply chain. 

Market Design, Forward Guidance, and Investment Decisions in Carbon Credit Markets (with Luis Alvarez, Victor Orestes, and Thiago Silva) \\
% We estimate the effects of forward guidance on the supply of carbon credits when trading is subject to over-the-counter (OTC) frictions, focusing on the CBIO market in Brazil. We combine the OTC tape data with firms' carbon credit holdings, balance sheet outcomes, and interfirm payments to study the impact on demand for carbon credits, borrowing, investment, and supply chain spillovers. We focus on the rapid increase in prices in June 2022 followed by a crash in July 2022, driven by speculation about forward guidance and an unexpected change in carbon credit policy. We show how low liquidity generated the volatility, and then propagated by limited float, insufficient hedging options, and the absence of designated market-makers. 


% TEACHING
\vspace{-5mm}
\section*{Teaching}
\textbf{CUHK Business School} \\
\teachingentry{Instructor, International Finance (FINA3020)}{Fall 2025}  \\[0.2em]
\textbf{MIT Economics} \\
\teachingentry{TA, Development Economics: Macroeconomic Issues (14.772)}{Spring 2025}
\teachingentry{TA, Digitization and Innovation in Developing Countries (14.193)}{Spring 2025}
\teachingentry{TA, Estimation and Inference for Linear Causal and Structural Models (14.381)}{Fall 2023}
\teachingentry{TA, Statistical Methods in Economics (14.380)}{Fall 2023}
\teachingentry{TA, Economics of Digitization (14.27)}{Spring 2022, Spring 2025}
\teachingentry{TA, The Challenge of World Poverty (14.73)}{Fall 2021}

\vspace{-5mm}
\section*{Fellowships and Grants}
\regentry{NSF Graduate Research Fellowship (138,000 USD total)}{2019–2024}
\regentry{George and Obie Shultz Fund (7,550 and 5,540 USD)}{2022, 2024}
\regentry{EU Horizon 2020 Research and Innovation Grant  (8,000 EUR)}{2022}
\regentry{STICERD (10,000 GBP)}{2022}
\regentry{PEDL Exploratory Research Grant (25,000 GBP)}{2022}

\section*{Previous Positions}
\entry{Pre-Doctoral Fellow and Research Assistant}{2016–2019}{University of Chicago, Professor Michael Greenstone}
\entry{Trading Intern}{2015}{Jane Street}

\section*{Presentations}
\subimport{../_generated/}{cv_presentations.tex}

\section*{Other}
Citizenship: U.S.\\
Languages: English (native), Mandarin Chinese (proficient), Spanish (intermediate), Portuguese (beginner)

\end{document}
